\documentclass[11pt,a4paper]{article}
\usepackage[T1]{fontenc}
\usepackage[utf8]{inputenc}
\usepackage{mathpazo}
\usepackage{microtype}
\usepackage[margin=1.1in]{geometry}
\usepackage{amsmath,amssymb,amsthm}
\usepackage{booktabs,tabularx,array}
\usepackage{enumitem}
\usepackage[most]{tcolorbox}
\usepackage{titlesec}
\usepackage[hidelinks]{hyperref}
\usepackage{setspace}
\setstretch{1.08}

\titleformat{\section}{\normalfont\Large\bfseries}{\thesection}{1em}{}
\titleformat{\subsection}{\normalfont\large\bfseries}{\thesubsection}{1em}{}

\theoremstyle{plain}
\newtheorem{proposition}{Proposition}
\theoremstyle{definition}
\newtheorem{definition}{Definition}
\theoremstyle{remark}
\newtheorem{remark}{Remark}

\newtcolorbox{caution}[1][]{colback=gray!6,colframe=gray!55,boxrule=0.5pt,arc=1pt,
  left=6pt,right=6pt,top=4pt,bottom=4pt,fonttitle=\bfseries\small,title={Caution#1}}
\newtcolorbox{statusbox}[1]{colback=white,colframe=black!70,boxrule=0.6pt,arc=0pt,
  left=8pt,right=8pt,top=6pt,bottom=6pt,fonttitle=\bfseries,title={#1}}

\newcommand{\tier}[1]{\textup{\textsc{[#1]}}}
\newcommand{\Eone}{\tier{E1}}
\newcommand{\M}{\mathcal{M}}
\title{\textbf{Load-Bearing Words}\\[0.4em]
\large Meaning, Consequence, Revision, and Deletion in Language}
\author{Flyxion\\ \small Independent Researcher}
\date{24 September 2026}
\begin{document}
\maketitle

\begin{abstract}
\noindent
Words are usually described as labels, and disputes about words as ``merely semantic.'' Yet a single word can end a career, collapse a clinical consensus, alter the reading of a peace settlement, or contribute to the outbreak of war. This essay proposes that a word's meaning is best understood not as what it names but as what it licenses: the set of responses, inferences, and actions it makes available to those who receive it. On this account a word is \emph{loaded} when different audiences derive sharply different licenses from it, and a word is \emph{load-bearing} when institutions have been built to act on its license. The essay applies the account to a series of cases: the clinical and colloquial lives of ``paranoia''; the contest over ``gender''; platform bans and the substitution vocabularies they produce; diplomatic texts in which a single article or a deleted clause changed a policy or started a war; and medical literature in which a short, decontextualized letter and an ambiguous abbreviation caused measurable harm. It then turns to deletion and revision, arguing that the ability to remove, retract, and correct words is as important to knowledge as the ability to produce them. Along the way it argues that load becomes measurable precisely where words bear institutional weight, since institutions supply the common response space the general notion lacks, and that chains of consequence contain damping stages as well as amplifying ones, whose placement is the central design question for consequential language. It traces the account to an older disagreement between Dryden, for whom words succeed by resemblance, and Schiller, for whom they succeed by the correctness of what they license. The analysis is set against David Deutsch's account of knowledge as error correction, with which it largely agrees, and against the Popperian advice not to quarrel about words, with which it partly disagrees: the advice is sound for inquiry and unsound wherever institutions allocate consequences by definition, and since research categories routinely flow into institutional use, the boundary between the two is better understood as a direction of flow than as a fixed line.
\end{abstract}

\begin{statusbox}{Statement of Scope}
\small
\begin{enumerate}[leftmargin=1.4em,itemsep=2pt]
\item This essay analyzes how words acquire consequence. It does not take a position on contested political or social questions, including those concerning gender, and presents disputed positions in terms their holders would recognize.
\item Historical cases are summarized from well-known accounts. Where historians disagree about the causal weight of a particular word, the disagreement is stated rather than resolved.
\item References were compiled during drafting and should be verified before citation.
\end{enumerate}
\end{statusbox}

\section{Introduction: The Myth of the Mere Word}

The phrase ``just semantics'' expresses a widespread intuition: that words are interchangeable containers for meanings that exist independently of them, so that arguments about words are arguments about packaging. The intuition is correct often enough to be useful. Two chemists who disagree about whether to call a compound an acid or a proton donor can usually resolve their difference by agreeing on the reaction.

It is wrong often enough to be dangerous. A patient's chart that reads ``10 U'' rather than ``10 units'' can be misread as 100. A telegram edited by deletion alone, with no word added, can turn a courteous royal refusal into a public insult. A resolution that omits the word ``the'' can sustain two incompatible readings for decades. A word posted in the wrong context can remove an account from a platform in seconds, while a word spoken in the right context, by the right person, into the right institution, can authorize the use of force.

None of these cases is ``merely semantic,'' and none is mysterious. Each can be explained by attending to what a word does rather than what it names. The purpose of this essay is to make that attention systematic.

\section{A Working Account of Meaning}

\subsection{Meaning as license}

The account adopted here draws on two familiar ideas. The first, associated with the later Wittgenstein, is that the meaning of a word is largely given by its use \cite{wittgenstein1953}. The second, from J.\,L.\ Austin, is that some utterances do not describe actions but perform them: a promise, a verdict, an order, a marriage vow \cite{austin1962}. Combining the two yields a single working definition.

\begin{definition}[License]
The \emph{license} of a word or phrase $w$, uttered in context $c$ and received by audience $a$, is the distribution $\M_a(w,c)$ over responses that $a$ takes to be warranted by receiving $w$ in $c$. Responses include inferences, beliefs, emotional reactions, and actions.
\end{definition}

On this view a word is meaningful because it changes what happens next. A word that changed no one's continuation would be meaningless in the only sense that matters for consequence.

\subsection{An older disagreement: resemblance and sign}

The license account has a longer history than its twentieth-century sources suggest, and the history is clearest where two earlier writers disagree.

In his preface to \emph{Annus Mirabilis} (1667), John Dryden describes poetic wit as a faculty of retrieval. Imagination, ``like a nimble spaniel, beats over and ranges through the field of memory,'' searching for ``the species or ideas'' of the things it wishes to represent, and good description succeeds when it ``sets before your eyes the absent object, as perfectly, and more delightfully than nature'' \cite{dryden1667}. On this model words work by resemblance. A sign is good when it reproduces its object, and the mind's stock of meanings is a stock of images.

Friedrich Schiller, writing a little over a century later in the voice of the young Julius in his \emph{Philosophical Letters} (1786), rejects the model outright. ``Our purest ideas are by no means images of things,'' Julius writes, ``but only their signs or symbols determined by necessity.'' Truth, accordingly, ``is no property of the idioms, but of the conclusion; it is not the likeness of the sign with the thing signified \ldots\ but the agreement of this conception with the laws of thought'' \cite{schiller1786}. His evidence is mathematics. Nothing in the letters $A$ and $B$ or the signs $+$ and $=$ resembles a comet, yet a calculation made in them brings the comet back on schedule, centuries later, to the corner of the sky where it was predicted.

Schiller's position is, in substance, the one adopted here. A sign earns its meaning not by looking like its object but by licensing correct continuations: the right prediction, the right inference, the right act. The comet is the limiting case of a word with no resemblance to its referent and complete reliability of license. Dryden's model, by contrast, cannot explain why an unpictorial word such as ``the'' in a Security Council resolution should matter at all, since there is nothing for it to resemble. The cases in this essay are, almost without exception, cases in which resemblance is irrelevant and license is everything.

\begin{remark}
The disagreement is not simply won by Schiller. Dryden is describing what makes description vivid, and vividness is itself a property with consequences: a vivid image propagates further and is acted on faster. Resemblance may be irrelevant to what a word licenses and still highly relevant to how strongly the license is felt. That distinction returns in the discussion of amplification below.
\end{remark}

\subsection{Loaded words}

\begin{definition}[Load]
A word is \emph{loaded} to the extent that different audiences receiving it in the same context derive different licenses from it. Schematically,
\[
L(w,c) \;=\; \mathbb{E}_{a,a'}\; D\big(\M_a(w,c)\,\big\|\,\M_{a'}(w,c)\big),
\]
where $D$ is any suitable divergence between response distributions and the expectation runs over pairs of audience members.
\end{definition}

A word with low load licenses roughly the same response in everyone. ``Seven'' is unloaded in most contexts. A word with high load licenses one response in one listener and a very different response in another, and the difference is typically not about reference. Two people may agree entirely about \emph{who} a word applies to and disagree completely about what follows from applying it.

\subsection{Load-bearing words}

\begin{definition}[Load-bearing]
A word is \emph{load-bearing} when an institution has been constructed to convert its license into action with little or no further deliberation.
\end{definition}

Load and load-bearing are independent. A slur may be heavily loaded and, in a private conversation, bear no institutional load at all; in a moderated forum it may bear the entire weight of an automated ban. A launch authentication code carries no load, since it means the same thing to every trained officer, and bears the greatest institutional load of any string of characters in existence. The most consequential cases are those in which the two properties coincide: a word that different audiences read differently, received by an institution that acts on one reading.

\begin{caution}
The formalism is schematic. It makes no claim that licenses can be measured precisely, only that the distinction between what a word names and what it licenses is what explains its consequences. The expression is used to keep that distinction in view, not to compute with it.
\end{caution}

\subsection{Where load becomes measurable}

An obvious objection is that responses of different kinds (inferences, beliefs, feelings, actions) share no common scale, so that a divergence between distributions over them is not well defined. The objection is correct in general, and it points to something more interesting than a defect.

Load-bearing institutions fix the response space themselves. A moderation system's responses to a post are a small finite set: leave it, label it, remove it, suspend the account. A pharmacist's responses to a prescription are a set of doses. A court's responses to a charge are verdicts and sentences. Once an institution has reduced the space of admissible continuations to such a set $R_I$, the license of a word \emph{for that institution} is a distribution over $R_I$, and divergence between the licenses different officials derive from the same word is an ordinary, estimable quantity. Inter-rater disagreement among content moderators, or variation among clinicians interpreting the same order, measures precisely this.

\begin{proposition}
Load is well defined, and in principle measurable, exactly where a word is load-bearing, because load-bearing institutions supply the common response space that the general definition lacks. Outside institutions, load remains a qualitative notion.
\end{proposition}

The result is not a coincidence. Institutions exist in large part to convert open-ended human responses into a bounded set of actions, and that conversion is what makes both their consequences and their disagreements countable.

\subsection{Amplification}

When a word bears institutional load, its license propagates. An order is passed down a hierarchy; a headline is republished; a citation is recited. If each stage multiplies consequences by a factor $g_i$, a single utterance produces an effect proportional to $\prod_i g_i$. When every $g_i > 1$, small differences in the initial word become large differences in outcome. Chains of this kind are what allow a sentence to contribute to a war.

Dryden's vividness belongs here. A description that ``sets before your eyes the absent object'' raises the gain of the stages that receive it: vivid accounts are repeated more readily and acted on more quickly than accurate but pallid ones. The same property that makes an image good poetry can make it dangerous evidence, which is why institutions that must act carefully often prefer flat language to vivid language.

Real chains also contain stages that shrink consequences rather than multiply them. A reviewer who asks a question, an officer who requests confirmation, a pilot who reads back a clearance, an appeal that pauses enforcement: each is a stage with gain below one. Writing $\mathcal{A}$ for amplifying stages and $\mathcal{D}$ for damping ones, the net effect of an utterance scales as
\[
G \;=\; \prod_{i\in\mathcal{A}} g_i \;\cdot\; \prod_{j\in\mathcal{D}} \delta_j, \qquad g_i > 1,\quad 0 \le \delta_j < 1 .
\]
A damping stage matters most when it is placed early, since it then acts on the error before amplification, and when its gain depends on the content it receives, passing a clear order nearly unchanged while stopping an ambiguous one almost completely. The read-back is the model case: it transmits a correct instruction at almost no cost and halts a misheard one. Several of the cases below are, on inspection, stories about the presence or absence of such a stage. Section~\ref{sec:engineer} argues that the placement of damping stages is the central design question for load-bearing language.
\section{Paranoia: A Diagnosis That Doubles as a Verdict}

\subsection{Two lives of one word}

``Paranoia'' comes from Greek roots meaning roughly ``beside the mind.'' In clinical usage it names a pattern of persistent, unwarranted suspicion, sometimes of delusional intensity, and it appears in the names of recognized conditions. In ordinary usage it names almost any suspicion the speaker considers excessive. The two usages share a word and very little else: the clinical term implies that the suspicion is unwarranted \emph{because} of a disorder; the colloquial term implies only that the speaker does not share it.

There is also a third life. In security, engineering, and management the word is used approvingly. A well-known business memoir made its title out of the claim that only the paranoid survive, and security practitioners speak of ``healthy paranoia'' as a professional virtue: the habit of assuming that a system can be attacked in ways not yet observed.

\subsection{Why the word is loaded}

The load of ``paranoia'' arises because a single token carries a descriptive component (``you suspect something'') and an adjudicative component (``your suspicion is baseless''), and the adjudicative component arrives without an argument. To call a concern paranoid is to dismiss it while appearing merely to describe it. The listener who hears the clinical overtone receives a license to discount the speaker entirely; the listener who hears the engineering sense receives a license to take the concern seriously. The same sentence licenses opposite continuations.

\begin{remark}
The effect has a structural cousin in terms such as ``conspiracy theorist'' and ``alarmist.'' Each can be used accurately. Each can also be used to convert an unexamined claim into an examined-and-rejected one without the examination taking place. The problem is not the word but the adjudication smuggled inside it.
\end{remark}

\subsection{Suspicion as a preserved alternative}

A companion essay, \emph{The Power of Paranoia}, argues that suspicion begins as a rational recognition that appearances underdetermine causes, and becomes destructive only when it is permitted to substitute for proof. On that account due process is disciplined paranoia: institutional machinery for keeping alternative explanations alive until evidence can decide among them. The present essay adds a linguistic corollary. Because the word ``paranoid'' performs adjudication, using it forecloses precisely the interval that due process exists to protect. It declares the alternative closed. A more exact vocabulary would separate the three things the word bundles: the \emph{presence} of suspicion, the \emph{evidence} for it, and the \emph{judgment} about that evidence.

\section{Gender: When Definition Allocates}

\subsection{Why this word, and why now}

``Gender'' is among the most contested words in contemporary English, and the contest is frequently described by each side as the other side's attempt to change the meaning of a word. Both descriptions contain something true, because the word is doing several jobs at once, and different communities have come to treat different jobs as primary.

In ordinary historical usage the word functioned largely as a polite synonym for sex. In mid-twentieth-century psychology and sociology it was given a technical sense distinguishing social roles and expectations from biological sex. In more recent usage it also names a person's internal sense of their own identity. A single token thus carries at least three candidate senses: a biological category, a social role, and a self-understanding.

\subsection{Why the dispute cannot be settled by dictionary}

If the dispute were only about labels, it could be resolved by stipulation, as the chemists resolved theirs. It cannot, because in many institutions the word is load-bearing. Definitions determine eligibility for sports categories, placement in single-sex facilities, entries on identity documents, eligibility for certain medical services, and data collection in research. Whoever fixes the sense of the word in a given institution thereby fixes an allocation.

Those who emphasize sex-based definitions generally argue that certain protections, categories, and data exist because of bodily sex, and that redefining the word dissolves those protections or corrupts the data. Those who emphasize identity-based definitions generally argue that social recognition and access to services should follow how people understand and live their own lives, and that sex-based definitions exclude or harm a vulnerable group. Each side sees the other's definition as an attempt to win an allocation dispute by fixing a word, and each is, in its own terms, correctly describing the stakes.

\begin{remark}
The medical literature illustrates why precision matters here without resolving the dispute. Some drugs are metabolized differently on average by female and male bodies; in 2013 United States regulators lowered the recommended starting dose of the sleep medication zolpidem for women for this reason. Research that records only one of the word's senses, or conflates them, loses information that the other sense would have preserved. Clinicians generally find both physiology and a patient's identity relevant to care, which is an argument for recording both clearly rather than for either definition prevailing.
\end{remark}

\subsection{What disaggregation looks like}

Disaggregation does not settle any of the underlying disputes. It changes their form, by replacing a contest over a token with explicit choices about criteria that can be argued on their merits. Three illustrations show the structure without favoring an outcome.

\begin{description}[style=nextline,leftmargin=1.2em,itemsep=4pt]
\item[Research instruments.] Some surveys and clinical research instruments ask separately about sex recorded at birth and about current gender identity, rather than asking a single question with one contested word. Each sense is then available for the analyses where it is relevant, and a later analyst can use either without reinterpreting the data.
\item[Identity documents.] A document field can name what it records: sex as recorded at birth, legal sex as currently recognized, or gender identity. Jurisdictions disagree, sometimes sharply, about which of these fields documents should carry and which should govern particular uses. Naming the field does not resolve that disagreement, but it prevents one field from being silently read as another.
\item[Sports eligibility.] A governing body can state the criterion it actually applies, for instance a physiological measure, a legal status, or a self-declared identity, rather than applying the word alone. Different bodies have adopted different criteria and defend them on grounds of fairness, safety, or inclusion. Once the criterion is explicit, argument can address whether it serves the category's purpose, which is the question actually in dispute.
\end{description}

In each case both sides can adopt the disaggregated structure and continue to disagree about which criterion should govern which use. What disaggregation removes is the specific failure mode in which an allocation is decided by whichever sense of a word an institution happens to have inherited, without anyone having argued for it.

\subsection{What the account adds}

On the license account, the contest over ``gender'' is a contest over which continuation an institution will attach to the word. That is why it is not ``just semantics,'' and also why disputants often talk past one another: they may agree on every fact about a particular person and disagree about which license the word should carry. Progress, where it occurs, tends to come from disaggregating the word, that is, from institutions stating which sense they are using for which purpose, rather than from one sense defeating the others.

\section{Banned Words and the Migration of Meaning}

\subsection{Filtering tokens}

Online platforms remove content and accounts partly on the basis of words. Keyword filters are cheap and fast, and some words, notably slurs, are strongly associated with harassment. But a filter operates on tokens, while harm operates on licenses. The mismatch produces two predictable failure modes.

The first is false positives. The best-known name for the phenomenon comes from an English town whose name contains an obscene substring and whose residents found themselves blocked by early filters. More consequentially, filters catch reporting, scholarship, quotation, reclaimed in-group usage, and discussion of the very harms they target. A historian quoting propaganda and a propagandist producing it may use identical words.

The second is substitution. When a token is filtered while its license is still wanted, speakers invent new tokens that carry the same license. Users of video platforms have developed euphemisms such as ``unalive'' for death and suicide, coined to evade automated suppression. The ban removes the word and leaves the meaning, which migrates to new vocabulary that the filter does not yet recognize. The filter must then chase the migration, and each round of chasing makes the vocabulary more opaque to outsiders, including moderators, researchers, and parents.

\subsection{Formal note}

On the license account, a ban deletes a token $w$ from the permitted vocabulary while leaving the target license $\M^\ast$ unchanged. Speakers then search for a substitute $w'$ such that $\M(w',c) \approx \M^\ast$. Because vocabulary is effectively unbounded, such a $w'$ can almost always be found. Token-level enforcement can raise the cost of expressing a license; it cannot remove the license, which lives in use rather than in spelling.

\begin{caution}
This is not an argument against moderation. Raising the cost of harassment has real value, and some words carry histories that justify strong restrictions. The argument is only that enforcement aimed at tokens will be outpaced by speakers aimed at licenses, and that context-sensitive judgment, though slower and costlier, is the only enforcement that tracks the thing that matters.
\end{caution}
\section{Words That Start and End Wars}

\subsection{Deletion as provocation: the Ems Dispatch}

In July 1870 King Wilhelm I of Prussia, taking the waters at Bad Ems, was approached by the French ambassador with a demand for assurances concerning the Spanish succession. The king declined, courteously, and sent an account of the exchange to his chancellor, Otto von Bismarck. Bismarck prepared a shortened version for the press. By his own later account he added nothing; he only cut. The condensed text made the king's refusal appear curt and the ambassador appear dismissed. Published in a charged atmosphere, it was received in Paris as an insult, and within days France declared war.

The case is the purest available demonstration that deletion is an act of authorship. Every sentence in the published text was, in some sense, true. What changed was the license: the removal of mitigating context converted a diplomatic exchange into a public humiliation, and two institutions, a press and a parliament, were ready to act on the humiliation.

\subsection{One word, two peaces: Resolution 242}

United Nations Security Council Resolution 242, adopted after the 1967 Arab--Israeli war, called for withdrawal ``from territories occupied in the recent conflict.'' The English text lacks a definite article. The French text, equally authoritative, reads in a way that many take to mean ``the territories.'' Whether the resolution required withdrawal from all occupied territory or from some of it has been argued ever since, with the absence of a single word cited on one side and the French text and the resolution's other provisions cited on the other. Whatever one concludes, the case shows that a text negotiated to the level of individual articles can remain load-bearing for generations precisely because of what it omits. The ambiguity of the word \emph{article} is apt here: in diplomatic texts an article is a numbered provision, and in this case the dispute turned on the other kind, a single grammatical article absent from one provision. Dryden, who dismissed wordplay as ``the jingle of a more poor Paronomasia,'' might not have approved of the observation, but the pun is the case.

\subsection{Constructive ambiguity}

Ambiguity is not always a defect. Negotiators sometimes use deliberately ambiguous language so that parties who could not agree on a precise formulation can each accept a text they read differently, deferring the conflict to a later time when it may be more tractable. The 1998 Good Friday Agreement in Northern Ireland is frequently cited as relying in part on such language. In license terms, constructive ambiguity is the deliberate creation of load: a text is written so that different audiences derive different continuations, each acceptable to its own side. It can end a war. It can also leave a debt that comes due when an institution must finally act on one reading.

\subsection{Choosing which letter to answer}

During the Cuban Missile Crisis of October 1962, the Soviet leadership sent two letters in quick succession: a first, more conciliatory message and a second, harder one adding a demand concerning American missiles in Turkey. The American response publicly answered the first letter as though the second had not arrived, while the Turkish issue was addressed through a separate private assurance. The crisis ended. The episode shows the reverse of the Ems case: a deliberate choice about which words to treat as operative, made in order to keep available a continuation that the harder text would have closed. In the terms of Section~2, it is a damping stage inserted by the receiver. The Ems Dispatch is the same structure run in reverse: an editor positioned where damping could have occurred used the position to amplify.

\subsection{A contested case: \emph{mokusatsu}}

In late July 1945 the Japanese prime minister responded to the Potsdam Declaration using the word \emph{mokusatsu}, which can mean to withhold comment, to take no notice of, or to treat with silent contempt. It was widely reported in the Allied press as a rejection. The episode is sometimes presented as a mistranslation that caused the atomic bombings. Most historians regard that claim as greatly overstated, since the relevant decisions were driven by much larger considerations. The case is included because it illustrates something the stronger claim distorts: when an institution is already poised to act, a loaded word can be read as confirmation of what it was prepared to do.

\subsection{Words that dehumanize}

During the 1994 genocide in Rwanda, which unfolded amid a civil war, broadcasts on the radio station RTLM repeatedly referred to Tutsi as \emph{inyenzi}, cockroaches, and called on listeners to act. An international tribunal later convicted figures associated with the station. Dehumanizing vocabulary works by changing license: a word that places a group of people in the category of pests licenses responses that no audience would accept toward neighbors. It is the clearest case in which the words themselves were treated by a court as part of the crime.

\section{Precision in Medicine}

\subsection{A letter without its context}

In 1980 the \emph{New England Journal of Medicine} published a one-paragraph letter reporting that, among a large population of hospitalized patients who had received narcotic analgesics, few had developed documented addiction \cite{porter1980}. The letter was brief and made no claim about outpatients or long-term prescribing. A 2017 bibliometric analysis found that it had been cited hundreds of times, the large majority of citations using it as evidence that addiction was rare in patients treated with opioids, and most omitting that the patients had been hospitalized \cite{leung2017}. The authors argued that this pattern of citation contributed to the North American opioid crisis by helping to reassure prescribers.

The case is a precise inverse of the Ems Dispatch. There, deletion of context was performed deliberately by one editor. Here, it was performed diffusely by hundreds of citing authors, each of whom repeated the finding and dropped the qualifying phrase. A single additional clause, ``in hospitalized patients,'' preserved in each citation, would have changed the license of the finding substantially.

\subsection{Abbreviations that kill}

Medication errors caused by ambiguous notation are common enough that hospital accreditation bodies maintain lists of prohibited abbreviations \cite{jointcommission}. The letter ``U'' for units can be read as a zero, turning ten units into one hundred. A trailing zero, as in ``1.0\,mg,'' can be read as ten milligrams if the decimal point is missed; a missing leading zero, as in ``.5\,mg,'' can be read as five. Abbreviations for different drugs have been confused with one another. These are the smallest possible textual differences, a mark or a character, and they are load-bearing in the most literal way, since they are executed by a person with a syringe.

\subsection{Retraction and its limits}

When a paper is shown to be wrong or fraudulent, journals retract it. Retraction is a formal deletion of license: the paper remains visible, marked, but is no longer to be relied upon. Retracted papers nonetheless continue to be cited, often without acknowledgment of the retraction, for years afterward. The most prominent example in recent medicine is the 1998 paper that proposed a link between the MMR vaccine and autism, retracted in 2010, whose influence on vaccination rates long outlasted its formal withdrawal. The lesson is that deletion from the record is not deletion from circulation. A license, once granted and propagated, is harder to revoke than it was to issue.

\section{Deletion, Addition, and the Single-Word Edit}\label{sec:deletion}

\subsection{Edit distance is not meaning distance}

The cases above share a structure. Let $s$ be a text and $s'$ an edited version. Let $d(s,s')$ be the edit distance between them, the number of characters or words changed, and let $\Delta(s,s') = D(\M(s)\,\|\,\M(s'))$ be the change in license. The intuition behind ``just semantics'' is that $\Delta$ grows roughly with $d$, so that small edits have small effects. The cases show that no such bound holds:
\[
\sup_{s,s'}\ \frac{\Delta(s,s')}{d(s,s')} \;=\; \infty .
\]
A missing article, a deleted clause, a character read as a zero: each has small $d$ and large $\Delta$. Conversely, a document can be extensively rewritten, with large $d$, and license exactly the same responses. Meaning is not a smooth function of text.

\subsection{Addition changes scope}

Adding material changes meaning as readily as removing it. A qualifying clause narrows a claim's scope; a new section can reframe everything before it; a footnote can convert an assertion into a report of someone else's assertion. The Porter and Jick letter shows the cost of an omitted scope clause. The converse is also true: a text can be made misleading by the \emph{addition} of accurate but selective context, which shifts the license without any false statement being made.

\subsection{Why deletion matters}

A culture that can only add to its record cannot correct it. Error correction requires, at minimum, the ability to mark a claim as withdrawn, and often the ability to remove it. This is why retraction, correction, erratum, and redaction are core scholarly practices, and why legal systems increasingly recognize rights to have personal data erased. Deletion is not the opposite of knowledge; it is one of its instruments.

Deletion is also dangerous, for exactly the same reason. The Ems Dispatch was created by deletion. Removing records can destroy evidence that later inquiry needs, and a right to erase can become a means of concealment. The resolution is not to prefer addition or deletion in general but to require that consequential deletions be \emph{recorded}: that the fact of a removal, its author, and its reason remain inspectable even when the removed content does not. A retraction notice is the model. It deletes the license and preserves the history.

\section{Knowledge, Error Correction, and the Quarrel About Words}

\subsection{Deutsch's frame}

David Deutsch's account of knowledge, developed in \emph{The Beginning of Infinity} \cite{deutsch2011} and applied to artificial intelligence in a 2020 talk titled ``Knowledge Creation and Its Risks'' \cite{deutsch2020}, holds that knowledge grows by conjecture and criticism, and that the traditions sustaining it are traditions of error correction. In that talk he argued against attempts to constrain the thinking of a genuinely general intelligence, likening such constraint to slavery, and in his wider work he treats the suppression of criticism as among the gravest threats to progress.

The present essay agrees with the core of this account and applies it to words. Retraction, correction, erratum, and recorded deletion are error-correcting institutions for language. Token bans that chase vocabulary without engaging licenses are, in Deutsch's terms, attempts to prevent a conjecture from being expressed rather than to criticize it, and they fail for the reason his account predicts: expression finds another route.

\subsection{Schiller's erring reason}

Schiller's Julius anticipates this account from an unexpected direction. He grants that his own speculations may be ``wholly an invention,'' and argues that this does not make them worthless: it was ``the provision of the Wisest Spirit that the erring reason should also people the chaotic world of dreams, and make fruitful even the barren ground of contradiction,'' since ``every facility of the reason, even in error, increases its readiness to accept truth'' \cite{schiller1786}. Stripped of its theology, this is a claim that error is an input to knowledge rather than merely its failure, which is close to the core of Deutsch's position.

Julius's own illustration makes the point better than he intended. He cites Columbus as a man who trusted a calculation, sailed into unknown water to find the land his reckoning required, and was proved right. The historical Columbus was not right. He had substantially underestimated the size of the Earth, was seeking a western route to Asia, and was saved from the consequences of his error by a continent that appeared on no one's map. The example is thus a case not of calculation vindicated but of error made fruitful by an encounter with reality it did not anticipate. That is precisely the process Julius elsewhere describes, in which ``all tortuous deviations of the wandering reason at length strike into the straight road,'' and it is the process by which, on Deutsch's account, knowledge grows.

For the argument of this essay the lesson is specific. Error can be fruitful only if it remains available for correction: Columbus's miscalculation was productive because a landfall contradicted it and the contradiction was recorded. An error that is erased without trace, or never exposed to the test that would contradict it, teaches nothing. This is a further reason for the requirement, stated in Section~\ref{sec:deletion}, that consequential deletions leave a record.

\subsection{Where the Popperian advice fails}

Deutsch's intellectual tradition inherits from Karl Popper a characteristic impatience with disputes about the meanings of words. Popper advised that one should not be drawn into taking questions about words and their meanings seriously, since what matters is the truth of theories, not the definitions of terms \cite{popper1976}. For inquiry the advice is excellent. A physicist gains nothing by arguing whether an electron ``really'' is a particle; the theory does the work.

The advice fails wherever words are load-bearing. When an institution allocates a sports category, a medication dose, a territorial withdrawal, or a platform ban on the basis of a word's sense, the definition is no longer a convention in service of inquiry. It is the allocation. The quarrel about words, in such settings, is the quarrel about outcomes conducted in the only terms the institution recognizes. One may still wish that institutions were structured to allocate by explicit criteria rather than by contested terms, and that wish is a good design principle. It does not make the existing quarrels trivial.

\subsection{When inquiry becomes allocation}

The distinction just drawn has a weakness: many words live on both sides of it at once. ``Gender'' is contested within academic disciplines, where the Popperian advice would counsel indifference to definitions, and load-bearing in hospitals and sports bodies, where it cannot. Nor is the boundary between inquiry and allocation itself fixed or uncontested.

A partial resolution is to treat the boundary as a direction of flow rather than a line. Definitions adopted in inquiry do not stay there. Diagnostic manuals are research instruments that become billing codes and eligibility criteria. Survey categories become the categories in which populations are counted, funded, and served. Terms coined to classify data become the terms in which policies are written. A word crosses from inquiry into allocation at the point where some institution begins to act on the categories inquiry produced.

\begin{proposition}
The Popperian advice to disregard disputes about words is sound only while the definitions in question remain upstream of any allocation. Once an institution acts on research categories, the definitions used in that research acquire allocative consequences, and disputes about them are no longer merely verbal.
\end{proposition}

In medicine and the social sciences this condition is rarely satisfied for long, since their categories are adopted by institutions almost as soon as they are useful. The practical implication is modest but real. Researchers who coin or refine a category in such fields should anticipate its institutional afterlife, and should prefer disaggregated categories, whose senses can be separated later, over single terms whose senses can only be fought over. Where inquiry and allocation cannot be kept apart, the quarrel about words cannot be avoided, but it can be conducted over explicit criteria rather than over tokens.

\section{Engineering Words That Cannot Be Misread}\label{sec:engineer}

Where words bear the greatest institutional load, institutions have learned to strip them of load. Military and aviation radio uses a phonetic alphabet so that ``B'' and ``D'' cannot be confused, and requires that critical instructions be read back to the sender. Hospitals prohibit ambiguous abbreviations. Nuclear command systems require authenticated orders, fixed formats, and more than one person to act, so that no ordinary sentence, however forcefully spoken, can by itself release a weapon.

Each of these practices is a damping stage in the sense of Section~2, deliberately placed where amplification would otherwise be greatest: read-back before execution, a prohibited notation before dispensing, a second authenticated officer before launch. They share a principle that summarizes the argument of this essay:

\begin{statusbox}{Principle}
The more institutional load a word bears, the less interpretive load it should carry. The safest load-bearing words are those engineered so that every competent audience derives the same license from them, and so that a misreading is caught before it is executed.
\end{statusbox}

The principle applies in reverse to the dangerous cases. ``Paranoid,'' used to dismiss, is a loaded word bearing adjudicative load without procedure. A filtered token is a load-bearing word whose license the filter cannot see. A decontextualized citation is a loaded claim bearing clinical load. A contested definition in an allocating institution is the principle's failure in its most visible form. In each case the remedy has the same shape: reduce the load the word carries, or add the procedure that should accompany the load it bears.

\section{Conclusion}

Words are meaningful because they change what happens next. They are loaded when they change it differently for different listeners, and load-bearing when institutions act on them without further thought. Most disputes dismissed as ``just semantics'' are disputes about license, and most catastrophes attributed to a single word are catastrophes of amplification, in which an institution was ready to act and a word supplied the occasion.

Three practical conclusions follow. First, when a word adjudicates while appearing to describe, as ``paranoid'' often does, separate the description from the judgment and demand the argument for the judgment. Second, when a definition allocates, state the sense and the purpose explicitly rather than fighting over the token. Third, protect the ability to delete and correct, and require that consequential deletions leave a record. A text can be changed by one word, one clause, or one missing article. Knowledge survives that fact only because it can also be changed back.

\begin{thebibliography}{99}
\small
\bibitem{austin1962} Austin JL. \emph{How to Do Things with Words}. Oxford: Clarendon Press (1962).
\bibitem{deutsch2011} Deutsch D. \emph{The Beginning of Infinity: Explanations That Transform the World}. London: Allen Lane (2011).
\bibitem{deutsch2020} Deutsch D. Knowledge creation and its risks. Talk and discussion, Leverhulme Centre for the Future of Intelligence, University of Cambridge (2020).
\bibitem{dryden1667} Dryden J. An account of the ensuing poem, in a letter to the Honorable Sir Robert Howard. Preface to \emph{Annus Mirabilis: The Year of Wonders, 1666} (1667).
\bibitem{jointcommission} The Joint Commission. Official ``Do Not Use'' list of abbreviations. \emph{[Verify current version.]}
\bibitem{leung2017} Leung PTM, et al. A 1980 letter on the risk of opioid addiction. \emph{New England Journal of Medicine} (2017). \emph{[Verify.]}
\bibitem{popper1976} Popper K. \emph{Unended Quest: An Intellectual Autobiography}. London: Fontana (1976). \emph{[Verify passage.]}
\bibitem{porter1980} Porter J, Jick H. Addiction rare in patients treated with narcotics. \emph{New England Journal of Medicine} (1980).
\bibitem{schiller1786} Schiller F. \emph{Philosophische Briefe} (\emph{Philosophical Letters}), including ``Theosophie des Julius'' (1786). Quoted from a nineteenth-century English translation. \emph{[Verify translation and edition.]}
\bibitem{wittgenstein1953} Wittgenstein L. \emph{Philosophical Investigations}. Oxford: Blackwell (1953).
\end{thebibliography}
\end{document}
